\documentclass[11pt,a4paper]{article}
\usepackage[margin=1in]{geometry}
\usepackage{amsmath,amssymb,amsthm}
\usepackage{algorithm}
\usepackage{algpseudocode}
\usepackage{booktabs}
\usepackage{hyperref}

\theoremstyle{definition}
\newtheorem{definition}{\textbf{Definition}}[section]
\newtheorem{theorem}{\textbf{Theorem}}[section]
\newtheorem{lemma}[theorem]{\textbf{Lemma}}
\newtheorem{remark}{\textbf{Remark}}[section]

\title{\textbf{Advanced DPO Variants: SimPO, KTO, and Reference-Free Alignment}}
\author{\textbf{Team Scaibu} \\ \small Email: \texttt{jaydeep.9664920749@gmail.com} \\ \small \textbf{Architecture and Core Engine Team}}
\date{\textbf{October 2026}}

\begin{document}
\maketitle

\section{Definitions and Cognitive Mental Models}

\subsection{Formal Mathematical Definition}
\textbf{Definition (Reference-Free Objectives, Length Normalization, and Unpaired Alignment):}
While Direct Preference Optimization (DPO) significantly simplifies alignment, standard DPO exhibits three critical practical limitations: (1) It requires caching or running a frozen reference model $\pi_{\text{ref}}$, doubling memory footprint; (2) Unnormalized sequence log-probabilities induce length bias, rewarding longer responses simply because their log-likelihood decreases more slowly; and (3) It strictly requires paired preference data $(y_w, y_l)$, discarding ubiquitous unpaired binary feedback (e.g. user thumbs-up or thumbs-down).

To address these limitations, several prominent mathematical variants have been developed:
\begin{enumerate}
    \item \textbf{Simple Preference Optimization (SimPO)}: SimPO eliminates the reference model entirely and incorporates explicit length normalization directly into the implicit reward:
    {\boldmath
    \begin{equation}
    r_{\text{SimPO}}(x, y) = \frac{\beta}{|y|} \sum_{t=1}^{|y|} \ln \pi_\theta(y_t \mid x, y_{<t})
    \end{equation}
    }
    The policy is optimized using a target separation margin $\gamma > 0$:
    {\boldmath
    \begin{equation}
    \mathcal{L}_{\text{SimPO}}(\theta) = - \mathbb{E}_{(x, y_w, y_l)} \left[ \ln \sigma\left( r_{\text{SimPO}}(x, y_w) - r_{\text{SimPO}}(x, y_l) - \gamma \right) \right]
    \end{equation}
    }
    \item \textbf{Kahneman-Tversky Optimization (KTO)}: Inspired by prospect theory in behavioral economics, KTO operates directly on unpaired binary feedback (single responses labeled desirable $y \in \mathcal{Y}_{\text{des}}$ or undesirable $y \in \mathcal{Y}_{\text{und}}$) without requiring pairwise comparisons:
    {\boldmath
    \begin{equation}
    \mathcal{L}_{\text{KTO}}(\theta) = \mathbb{E}_{y \in \mathcal{Y}_{\text{des}}}\left[ 1 - \sigma\left( \beta \ln \frac{\pi_\theta(y)}{\pi_{\text{ref}}(y)} - z_{\text{ref}} \right) \right] + \lambda_D \mathbb{E}_{y \in \mathcal{Y}_{\text{und}}}\left[ 1 - \sigma\left( z_{\text{ref}} - \beta \ln \frac{\pi_\theta(y)}{\pi_{\text{ref}}(y)} \right) \right]
    \end{equation}
    }
    where $z_{\text{ref}}$ is a reference KL baseline and $\lambda_D > 1$ represents human loss aversion.
    \item \textbf{Identity Preference Optimization (IPO)}: Replaces the logistic loss with a quadratic penalty around target margin $\frac{1}{2\tau}$, eliminating gradient saturation on deterministic labels:
    {\boldmath
    \begin{equation}
    \mathcal{L}_{\text{IPO}}(\theta) = \mathbb{E}\left[ \left( \beta \ln \frac{\pi_\theta(y_w)}{\pi_{\text{ref}}(y_w)} - \beta \ln \frac{\pi_\theta(y_l)}{\pi_{\text{ref}}(y_l)} - \frac{1}{2\tau} \right)^2 \right]
    \end{equation}
    }
\end{enumerate}

\subsection{Expanded Non-Technical Definition (Intuition for Anyone)}
\textbf{Intuitive Conceptual Definition:}
\begin{enumerate}
    \item \textbf{SimPO (No Reference Model, No Verbosity Hack)}: DPO had a secret flaw: it loved long answers. SimPO fixes this by averaging the score per word ($1/|y|$). A short, punchy, 20-word answer that gets straight to the point is rewarded far more than a 200-word rambling essay! Furthermore, it throws away the reference model entirely.
    \item \textbf{KTO (Thumbs-Up / Thumbs-Down)}: Most real-world products (like Netflix, YouTube, or customer support bots) do not ask users to choose between two parallel answers. Users simply click ``Thumbs Up'' or ``Thumbs Down''. KTO enables models to learn directly from isolated thumbs-up and thumbs-down signals!
    \item \textbf{Loss Aversion in AI}: In human psychology (Kahneman-Tversky), losing \$10 feels twice as painful as winning \$10 feels good. KTO incorporates this asymmetry ($\lambda_D$): bad answers are penalized much more aggressively than good answers are rewarded, preventing toxic hallucinations.
\end{enumerate}

\subsection{Child-Friendly Mental Model (Physical Metaphor)}
Imagine a school grading homework:
\begin{itemize}
    \item \textbf{The Heavy Essay Problem (Standard DPO)}: A student submits a 50-page essay full of fluff and gets an A simply because the teacher gets tired of reading.
    \item \textbf{The SimPO Teacher}: Grades essays based on density of good ideas per page ($1/|y|$). If you explain the entire solar system in 3 brilliant paragraphs, you get the gold star!
    \item \textbf{The KTO Sticker System}: The teacher does not compare two kids' papers side-by-side. If a paper is good, it gets a gold star. If a paper has mean words, it gets a red sticker. The class quickly learns to earn gold stars and avoid red stickers!
\end{itemize}

\section{Core Mathematical Equations and Definitions}

\subsection{SimPO Length-Normalized Implicit Reward}
{\boldmath
\begin{equation}
r_{\text{SimPO}}(x, y) = \frac{\beta}{|y|} \ln \pi_\theta(y \mid x) = \frac{\beta}{|y|} \sum_{t=1}^{|y|} \ln \pi_\theta(y_t \mid x, y_{<t})
\end{equation}
}
\begin{itemize}
    \item \textbf{Formal Definition}: Average per-token log-likelihood scaled by regularization constant $\beta$.
    \item \textbf{What It Means}: Measures the average confidence per token, stripping away artificial length accumulation.
    \item \textbf{Why It Is Important}: Completely eliminates length bias without needing length penalty calibration or hyperparameter gymnastics.
\end{itemize}

\subsection{SimPO Target Separation Objective}
{\boldmath
\begin{equation}
\mathcal{L}_{\text{SimPO}} = - \frac{1}{B}\sum_{i=1}^B \ln \sigma\left( r_{\text{SimPO}}(x_i, y_{w, i}) - r_{\text{SimPO}}(x_i, y_{l, i}) - \gamma \right)
\end{equation}
}
\begin{itemize}
    \item \textbf{Formal Definition}: Logistic loss enforcing an explicit target separation margin $\gamma > 0$ between candidate rewards.
    \item \textbf{What It Means}: Penalizes the model unless the preferred response beats the rejected response by at least margin $\gamma$.
    \item \textbf{Why It Is Important}: Yields state-of-the-art alignment benchmarks while cutting training memory by 50\% (no reference model).
\end{itemize}

\subsection{KTO Prospect Theory Value Function}
{\boldmath
\begin{equation}
v(z) = \begin{cases} 1 - \sigma(z - z_{\text{ref}}) & \text{if desirable } (z > 0) \\ \lambda_D (1 - \sigma(z_{\text{ref}} - z)) & \text{if undesirable } (z < 0) \end{cases}
\end{equation}
}
\begin{itemize}
    \item \textbf{Formal Definition}: Asymmetric S-shaped utility function reflecting human loss aversion on unpaired feedback.
    \item \textbf{What It Means}: Scales gradient updates for undesirable examples by loss aversion parameter $\lambda_D > 1$.
    \item \textbf{Why It Is Important}: Allows training directly on unstructured production telemetry (single user upvotes and downvotes).
\end{itemize}

\section{Rigorous Mathematical Analysis Checklist}

\begin{itemize}
    \item \textbf{Notation}: $\beta$ (reward multiplier), $\gamma$ (SimPO target margin), $|y|$ (token length), $\lambda_D$ (KTO loss aversion parameter).
    \item \textbf{Epistemic Status}: SimPO developed by Meng et al. (Princeton, NeurIPS 2024); KTO developed by Ethayarajh et al. (Contextual AI, ICML 2024).
    \item \textbf{Dependency Graph}: Dataset $\to (\text{Length Normalization} \mid \text{Unpaired Flag}) \to \text{Variant Objective} \to \nabla_\theta \mathcal{L}$.
    \item \textbf{Dimensions}: Batch size $B$, token sequence lengths $|y_w|, |y_l|$.
    \item \textbf{Regimes}: Reference-free alignment, compact resource training, unpaired user telemetry fine-tuning.
    \item \textbf{Invariants}: Under SimPO, if chosen and rejected responses have identical length-normalized log-probs, the reward margin is identically $0$.
    \item \textbf{Isomorphisms}: SimPO is isomorphic to margin ranking on average cross-entropy densities.
\end{itemize}

\section{Theoretical Theorems and Formal Proofs}

\begin{theorem}[Invariance of SimPO Reward to Repetitive Token Stuttering]
Let $y$ be a token sequence of length $L$ with average log-likelihood $\bar{\ell} = \frac{1}{L}\sum_{t=1}^L \ln \pi(y_t)$. If the sequence is concatenated $k$ times to produce an identical stuttered sequence $y^{(k)}$ of length $kL$, the SimPO reward satisfies $r_{\text{SimPO}}(y^{(k)}) \equiv r_{\text{SimPO}}(y)$, whereas standard DPO implicit reward scales linearly: $r_{\text{DPO}}(y^{(k)}) = k \cdot r_{\text{DPO}}(y)$.
\end{theorem}

\begin{proof}
By definition of the stuttered sequence, each token block is repeated $k$ times:
{\boldmath
\begin{equation}
|y^{(k)}| = k \cdot L, \quad \sum_{t=1}^{kL} \ln \pi(y_t^{(k)}) = k \sum_{t=1}^L \ln \pi(y_t)
\end{equation}
}
Evaluating the SimPO implicit reward:
{\boldmath
\begin{align}
r_{\text{SimPO}}(y^{(k)}) &= \frac{\beta}{|y^{(k)}|} \sum_{t=1}^{kL} \ln \pi(y_t^{(k)}) = \frac{\beta}{kL} \left( k \sum_{t=1}^L \ln \pi(y_t) \right) \\
&= \frac{\beta}{L} \sum_{t=1}^L \ln \pi(y_t) = r_{\text{SimPO}}(y)
\end{align}
}
Now evaluate the standard DPO unnormalized implicit reward (assuming reference log-prob also scales by $k$):
{\boldmath
\begin{equation}
r_{\text{DPO}}(y^{(k)}) = \beta \sum_{t=1}^{kL} \left( \ln \pi(y_t^{(k)}) - \ln \pi_{\text{ref}}(y_t^{(k)}) \right) = k \cdot r_{\text{DPO}}(y)
\end{equation}
}
If $r_{\text{DPO}}(y) > 0$, multiplying by $k$ inflates the reward arbitrarily without improving quality. SimPO's length normalization cancels $k$ identically, proving mathematical immunity to repetitive length inflation.
\end{proof}

\begin{theorem}[Convexity and Non-Saturation of the IPO Objective]
The Identity Preference Optimization loss $\mathcal{L}_{\text{IPO}}(h) = (h - c)^2$ is strongly convex with non-vanishing gradients everywhere except at the exact target margin $h = c$.
\end{theorem}

\begin{proof}
Let $h$ represent the implicit log-ratio difference and $c = \frac{1}{2\tau}$.
The objective is $\ell(h) = (h - c)^2$.
We compute the first and second derivatives:
{\boldmath
\begin{align}
\ell'(h) &= 2(h - c) \\
\ell''(h) &= 2 > 0, \quad \forall h \in \mathbb{R}
\end{align}
}
Because $\ell''(h) \ge 2 > 0$ uniformly on $\mathbb{R}$, $\ell(h)$ is strongly convex with parameter $\mu = 2$.
Unlike the logistic sigmoid loss $-\ln \sigma(h)$ where the gradient $\sigma(-h) \to 0$ as $h \to \infty$, the derivative $|\ell'(h)| = 2|h - c|$ increases linearly with deviation from the target margin $c$.
Thus, IPO prevents gradient saturation on deterministic labels.
\end{proof}

\section{Foundational Architectural Questions and Epistemic Meta-Questions}

\subsection{Architectural Question 1: Eliminating the Reference Model}
\textbf{Question:} Why is eliminating the reference model in SimPO advantageous for multi-node training clusters?\\
\textbf{Answer:} In large distributed clusters (e.g. 70B+ parameter models across 64 GPUs), retaining a reference model requires either dedicating 50\% of GPU memory to frozen parameters or synchronizing activations across dedicated nodes. SimPO frees 50\% of accelerator memory, allowing $2\times$ larger batch sizes, higher sequence lengths, and eliminating cross-network reference communication.

\textbf{Meta-Question:} Can a model trained with SimPO drift away from grammatical fluency without a reference model?\\
\textbf{Answer:} Yes, if $\gamma$ is set too aggressively. Setting the target margin $\gamma \in [0.5, 1.4]$ provides an optimal balance: it guides preference separation while keeping the model anchored to high-likelihood generation regions.

\subsection{Architectural Question 2: Unpaired Feedback in Production}
\textbf{Question:} When is KTO strictly necessary over paired DPO variants?\\
\textbf{Answer:} When operating in real-world software applications where end-users interact asynchronously. Generating two candidate answers per query for human comparison is $2\times$ more expensive and delays user responses. KTO allows an application to serve single fast answers, collect passive thumbs-up/down feedback, and fine-tune models directly on the natural feedback stream.

\textbf{Meta-Question:} How does KTO balance positive versus negative feedback ratios?\\
\textbf{Answer:} In production telemetry, positive feedback often exceeds negative feedback by $10:1$ (or vice-versa). KTO incorporates the loss aversion weight $\lambda_D$ and reference baseline $z_{\text{ref}}$ to maintain invariant gradient calibration regardless of class imbalance in the feedback buffer.

\section{Algorithmic Implementation Specification}

\begin{algorithm}[H]
\caption{Unified DPO Variant Loss Evaluation (SimPO, KTO, IPO)}
\begin{algorithmic}[1]
\Require Log-probs $\mathbf{lp}_w, \mathbf{lp}_l$, lengths $\mathbf{L}_w, \mathbf{L}_l$, mode $type$, reference log-probs, parameters $\beta, \gamma, \lambda_D$
\Ensure Variant loss $\mathcal{L}_{\text{var}}$, margin scores $\mathbf{S}$, accuracy $\operatorname{Acc}$
\State Initialize $loss\_accum \gets 0.0, \quad correct \gets 0$
\State Allocate $\mathbf{S} \in \mathbb{R}^B$
\For{$i = 0$ \textbf{to} $B-1$}
    \If{$type == \text{"simpo"}$}
        \State $rw \gets (\beta / \mathbf{L}_w[i]) \cdot \mathbf{lp}_w[i]$
        \State $rl \gets (\beta / \mathbf{L}_l[i]) \cdot \mathbf{lp}_l[i]$
        \State $m \gets rw - rl, \quad \mathbf{S}[i] \gets m$
        \If{$m > \gamma$} \State $correct \gets correct + 1$ \EndIf
        \State $z \gets m - \gamma$
        \State $loss\_accum \gets loss\_accum + (\ln(1 + \exp(-z)) \text{ if } z \ge 0 \text{ else } -z + \ln(1 + \exp(z)))$
    \ElsIf{$type == \text{"ipo"}$}
        \State $diff \gets (\mathbf{lp}_w[i] - \mathbf{ref}_w[i]) - (\mathbf{lp}_l[i] - \mathbf{ref}_l[i])$
        \State $h \gets \beta \cdot diff, \quad \mathbf{S}[i] \gets h$
        \If{$h > 0.0$} \State $correct \gets correct + 1$ \EndIf
        \State $loss\_accum \gets loss\_accum + (h - 1.0 / (2.0 \cdot \beta))^2$
    \ElsIf{$type == \text{"kto"}$}
        \State $ratio \gets \beta \cdot (\mathbf{lp}_w[i] - \mathbf{ref}_w[i]), \quad \mathbf{S}[i] \gets ratio$
        \If{$is\_good[i]$}
            \State $prob \gets 1.0 / (1.0 + \exp(-\operatorname{clamp}(ratio, -50, 50)))$
            \State $loss\_accum \gets loss\_accum + (1.0 - prob)$
            \If{$ratio > 0.0$} \State $correct \gets correct + 1$ \EndIf
        \Else
            \State $prob \gets 1.0 / (1.0 + \exp(-\operatorname{clamp}(-ratio, -50, 50)))$
            \State $loss\_accum \gets loss\_accum + (1.0 - prob)$
            \If{$ratio < 0.0$} \State $correct \gets correct + 1$ \EndIf
        \EndIf
    \EndIf
\EndFor
\State $\mathcal{L}_{\text{var}} \gets loss\_accum / B, \quad \operatorname{Acc} \gets correct / B$
\State \Return $(\mathcal{L}_{\text{var}}, \mathbf{S}, \operatorname{Acc})$
\end{algorithmic}
\end{algorithm}

\section{Big-$\Theta$ Complexity Analysis}

\begin{itemize}
    \item \textbf{Length-Normalized Reward Evaluation Time}: $\Theta(B)$ division and scalar operations.
    \item \textbf{Separation Margin Loss Time}: $\Theta(B)$ numerically stable softplus steps.
    \item \textbf{Unpaired KTO Utility Time}: $\Theta(B)$ logistic evaluations.
    \item \textbf{Space Complexity}: $\Theta(B)$ memory buffer for margin tracking.
    \item \textbf{Parallel Reduction Depth}: $\mathcal{D} = \Theta(\log B)$ for batch loss reduction.
\end{itemize}

\section{Single Canonical Repository Reference}

The complete deterministic scalar implementation, verification suite, and unit test benchmarks are published at:
\begin{center}
\url{https://github.com/Chief-Strategist-J/llm-obs-infra}
\end{center}

\end{document}
